% TOP_PROBLEM: {"id": 7500118, "title": "Strength of the Auslander–Ellis Theorem", "category_id": 18, "category": {"id": 18, "name": "logic", "display_name": "Logic", "description": "Problems in mathematical logic, model theory, proof theory, and finite model theory.", "slug": "logic", "order_index": 18, "created_at": "2026-07-31T15:26:25.671Z"}, "status": "partially_solved"}
% FIRST_POSTED: null
% ENTRY_KIND: statement_only
% Imported from: batch20.tex; UnsolvedMath ID 7500118
\documentclass[11pt]{article}
\usepackage[a4paper,margin=25mm]{geometry}
\usepackage{amsmath,amssymb,mathtools}
\usepackage{fontspec,unicode-math}
\setmainfont{Latin Modern Roman}
\setsansfont{Latin Modern Sans}
\setmathfont{Latin Modern Math}
\usepackage{xurl,hyperref}
\hypersetup{colorlinks=true,linkcolor=blue,urlcolor=blue,pdftitle={UnsolvedMath Problem 7500118}}
\newcommand{\sref}[2]{\hyperlink{source:#1:#2}{[S#2]}}
\newcommand{\cataloguescope}[1]{\paragraph{Mathematical question.}#1}
\begin{document}
% UnsolvedMath ID 7500118; source catalogue code AMR-074-0118
\setcounter{section}{7500117}
\section{Strength of the Auslander–Ellis Theorem}\label{7500118}
{\small\sffamily UnsolvedMath ID 7500118 · Logic}
\subsection{Definitions and mathematical statement}

\paragraph{Shared notation.}$\mathbb N=\{1,2,\ldots\}$, and $\mathbb Z,\mathbb Q,\mathbb R,\mathbb C$ denote the integers, rationals, reals and complex numbers. Any other conventions are specified locally.


Work in second-order arithmetic over $\mathsf{RCA}_0$, the base theory with $\Delta^0_1$ comprehension and $\Sigma^0_1$ induction. The strength of a principle means which subsystems prove it and which subsystems it implies over that base; equivalence requires both implications.  Here $\mathsf{ACA}_0$ adds comprehension for all arithmetical formulas. For a set $X$, its Turing jump $X'$ codes which oracle computations with oracle $X$ halt; $X^{(\omega)}=\bigoplus_{n\geq0}X^{(n)}$ joins all finite jumps. The theory $\mathsf{ACA}_0^+$ is $\mathsf{RCA}_0+\forall X\,\exists X^{(\omega)}$. The Auslander--Ellis theorem says that for every compact metric space $(X,d)$, continuous $T:X\to X$ and $x\in X$, there is $y\in X$ which is uniformly recurrent and proximal to $x$. Uniform recurrence means $\forall\varepsilon>0\,\exists m\geq1\,\forall n\geq0\,\exists k<m:\ d(T^{n+k}y,y)<\varepsilon$. Proximality means that, for every $\varepsilon>0$, infinitely many $n$ satisfy $d(T^nx,T^ny)<\varepsilon$.
\paragraph{Statement recorded in the supplied source.}
Is the Auslander--Ellis theorem equivalent to $\mathsf{ACA}_0$ over $\mathsf{RCA}_0$? \sref{7500118}{2}
\subsection{Short English statement}
Does the existence of a uniformly recurrent point proximal to each starting point have exactly the strength of arithmetic comprehension?
\subsection{Sources}
\begin{description}
\item[\hypertarget{source:7500118:1}{S1}] ULAM.AI and the UnsolvedMath Contributors, UnsolvedMath: A Curated Collection of Open Mathematics Problems, record 7500118 (AMR-074-0118). CC BY 4.0. \url{https://huggingface.co/datasets/ulamai/UnsolvedMath}.
\item[\hypertarget{source:7500118:2}{S2}] Antonio Montalbán, Open Questions in Reverse Mathematics, author preprint dated 17 August 2010, Question 18, printed p. 7--8. \url{https://math.berkeley.edu/~antonio/papers/questionsRM.pdf}.
\end{description}
\label{end:7500118}
\end{document}
